%═══════════════════════════════════════════════════════════════════════════════
\chapter{Optimización iterativa: del descenso de gradiente a los métodos adaptativos}
\label{ch:cap5}
%═══════════════════════════════════════════════════════════════════════════════

El Capítulo~\ref{ch:cap4} construyó el MLP como objeto matemático preciso y estableció
la notación del forward pass. El Capítulo~\ref{ch:cap3} mostró que el paisaje de
$\mathcal{L}(\bm{\theta})$ es no convexo y que no existe forma cerrada para
$\bm{\theta}^* = \argmin\mathcal{L}$. El presente capítulo cierra ese hueco: diseña
y analiza los algoritmos iterativos que ajustan $\bm{\theta}$.

La perspectiva organizadora del capítulo es la siguiente: los algoritmos de
optimización son integradores numéricos de ecuaciones diferenciales. El descenso
de gradiente es el método de Euler aplicado a la ODE del flujo de gradiente. El
SGD introduce un término estocástico que convierte esa ODE en una SDE de Langevin
cuya dinámica obedece una ecuación de Fokker-Planck. Los optimizadores adaptativos
son integradores con paso variable que aproximan la acción del hessiano inverso.
Esta perspectiva no es una metáfora: produce cotas de convergencia precisas,
explica el rol de la tasa de aprendizaje como temperatura termodinámica, y conecta
directamente con las herramientas del cálculo estocástico que el físico y el
matemático aplicado necesitarán en el Capítulo~\ref{ch:cap12} para los modelos
generativos.

%───────────────────────────────────────────────────────────────────────────────
\section{El problema de optimización iterativa}
\label{sec:optim_iterativa}
%───────────────────────────────────────────────────────────────────────────────

\dificultad{2}{2}

\subsection{Por qué no hay solución de forma cerrada}
\label{subsec:no_forma_cerrada}

La condición necesaria de primer orden para un mínimo de $\mathcal{L}$ es
$\nabla_{\bm{\theta}}\mathcal{L}(\bm{\theta}^*) = \bm{0}$. Para la familia
lineal, esta ecuación es el sistema lineal $\bm{V}^\top\bm{V}\bm{\theta} =
\bm{V}^\top\bm{y}$ con solución exacta $\bm{\theta}^* = \bm{V}^\dagger\bm{y}$
(Capítulo~\ref{ch:cap2}). Para el MLP, la ecuación $\nabla_{\bm{\theta}}\mathcal{L}
(\bm{\theta}^*) = \bm{0}$ es un sistema de $p$ ecuaciones no lineales en $p$
incógnitas sin estructura algebraica explotable. No existe, en general, ningún
método algebraico para resolverlo.

El método de Newton resolvería el sistema iterativamente: partiendo de $\bm{\theta}^{(0)}$,
cada iteración resuelve $\bm{H}(\bm{\theta}^{(k)})\bm{d}^{(k)} = -\nabla\mathcal{L}
(\bm{\theta}^{(k)})$ y actualiza $\bm{\theta}^{(k+1)} = \bm{\theta}^{(k)} +
\bm{d}^{(k)}$. El problema es que $\bm{H}(\bm{\theta}) \in \mathbb{R}^{p\times p}$:
para un MLP con $p = 10^9$ parámetros, almacenar $\bm{H}$ requeriría
$\sim 10^{18}$ bytes (un exabyte), y factorizarla costaría $O(p^3) \approx 10^{27}$
operaciones. Newton es computacionalmente inviable a la escala de las redes
neuronales modernas.

Se necesita un algoritmo que use solo información de primer orden
($\nabla\mathcal{L}$, cuyo cómputo se desarrollará en el Capítulo~\ref{ch:cap6}),
requiera memoria $O(p)$, y escale linealmente con $p$.

\subsection{La idea central: seguir el gradiente negativo}
\label{subsec:idea_gradiente}

La Proposición~\ref{prop:max_ascenso} estableció que $-\nabla_{\bm{\theta}}\mathcal{L}
(\bm{\theta})$ es la dirección de máximo descenso local de $\mathcal{L}$ en
$\bm{\theta}$. El algoritmo más simple que explota esta propiedad es:
\begin{equation}
  \label{eq:gd}
  \bm{\theta}^{(k+1)} = \bm{\theta}^{(k)} - \eta\,\nabla_{\bm{\theta}}\mathcal{L}
  (\bm{\theta}^{(k)}),
\end{equation}
donde $\eta > 0$ es la \textbf{tasa de aprendizaje} (\textit{learning rate}).
Esta regla de actualización se denomina \textbf{descenso de gradiente} (GD,
\textit{gradient descent}). A lo largo de este capítulo se irá enriqueciendo con
mecanismos adicionales (estocasticidad, momento, adaptación), pero el núcleo
permanece: moverse en la dirección opuesta al gradiente.

%───────────────────────────────────────────────────────────────────────────────
\section{Descenso de gradiente: el método de Euler sobre el flujo de gradiente}
\label{sec:gd}
%───────────────────────────────────────────────────────────────────────────────

\dificultad{2}{3}

\subsection{La ODE del flujo de gradiente}
\label{subsec:flujo}

El descenso de gradiente en tiempo continuo es la ecuación diferencial ordinaria:
\begin{equation}
  \label{eq:flujo_gradiente}
  \dot{\bm{\theta}}(t) = -\nabla_{\bm{\theta}}\mathcal{L}(\bm{\theta}(t)),
  \qquad \bm{\theta}(0) = \bm{\theta}^{(0)}.
\end{equation}
Esta ecuación describe la trayectoria de un punto que sigue en todo momento la
dirección de máximo descenso de $\mathcal{L}$. Recibe el nombre de \textbf{flujo
de gradiente} de $\mathcal{L}$.

\begin{observacion}{Analogía con la mecánica}
La ecuación~\eqref{eq:flujo_gradiente} es la ecuación de movimiento de una
partícula en el potencial $\mathcal{L}(\bm{\theta})$ con amortiguamiento
infinito (sin inercia): la fuerza $-\nabla\mathcal{L}$ se equilibra instantáneamente
con la fricción. En mecánica, esta es la dinámica de una partícula sobre-amortiguada
o la ecuación de difusión de Smoluchowski sin ruido térmico. A diferencia de la
ecuación de Newton $\ddot{\bm{\theta}} = -\nabla\mathcal{L}$, que puede oscilar
alrededor de un mínimo, el flujo de gradiente descende monótonamente: $\frac{d}{dt}
\mathcal{L}(\bm{\theta}(t)) = \langle\nabla\mathcal{L}, \dot{\bm{\theta}}\rangle
= -\|\nabla\mathcal{L}\|_2^2 \leq 0$.
\end{observacion}

Para funciones convexas, la solución de~\eqref{eq:flujo_gradiente} converge al
mínimo global. Para funciones no convexas, converge a un punto crítico cuya
naturaleza depende de la inicialización.

\begin{fisicaconexion}{Descenso de gradiente como dinámica viscosa}
El flujo de gradiente $\dot{\bm{\theta}} = -\nabla\mathcal{L}(\bm{\theta})$ es
la ecuación de movimiento de una partícula en un campo de fuerzas conservativo
$\bm{F} = -\nabla\mathcal{L}$ con amortiguamiento viscoso infinito: toda la
energía cinética se disipa instantáneamente. Esta es exactamente la
\textbf{ecuación de Smoluchowski} de la física de coloides en el límite de
alta viscosidad (número de Reynolds $\to 0$). En ese límite, $m\ddot{\bm{r}} \approx 0$
y la ecuación de Newton se reduce a $\gamma\dot{\bm{r}} = \bm{F}$, con $\gamma$
el coeficiente de fricción. La descarga de un condensador a través de una
resistencia ($V = IR$, $V = Q/C$, $\dot{Q} = -Q/RC$) sigue la misma ecuación:
$Q(t) = Q_0 e^{-t/RC}$, donde $RC$ es la constante de tiempo análoga a $1/\mu_{\min}$
(el modo más lento del flujo de gradiente). La condición de estabilidad de Euler
para el condensador es $\eta < 2RC$: exactamente la misma cota que para el
descenso de gradiente.
\end{fisicaconexion}

\subsection{El descenso de gradiente como discretización de Euler}
\label{subsec:euler}

\dificultad{2}{3}

El método de Euler con paso temporal $\eta$ aplicado a~\eqref{eq:flujo_gradiente}
produce:
\[
  \bm{\theta}(t+\eta) \approx \bm{\theta}(t) + \eta\,\dot{\bm{\theta}}(t)
  = \bm{\theta}(t) - \eta\,\nabla_{\bm{\theta}}\mathcal{L}(\bm{\theta}(t)),
\]
que es exactamente la actualización~\eqref{eq:gd} con $k$ como índice de tiempo
discreto. La tasa de aprendizaje $\eta$ es el paso temporal del integrador.

Esta identificación tiene consecuencias inmediatas. La condición de estabilidad
del método de Euler para una ODE lineal $\dot{u} = \lambda u$ es $|\eta\lambda| < 2$
(región de estabilidad del plano complejo). Aplicada al flujo de gradiente
linealizado alrededor de un punto crítico $\bm{\theta}^*$:
$\delta\dot{\bm{\theta}} = -\bm{H}(\bm{\theta}^*)\delta\bm{\theta}$, los
autovalores de la dinámica son $\lambda_k = -\mu_k$ (los autovalores del hessiano,
con signo negativo). La condición de estabilidad de Euler requiere
$\eta|\mu_k| < 2$ para todo $k$, lo que da $\eta < 2/\mu_{\max}$: exactamente
la condición de la Proposición~\ref{prop:conv_cuad}. Las propiedades de
convergencia del descenso de gradiente son propiedades de estabilidad del
integrador de Euler para la ODE linealizada.

\subsection{Convergencia para funciones fuertemente convexas}
\label{subsec:conv_fuerte}

\dificultad{2}{2}

\begin{definicion}{Función fuertemente convexa}
\label{def:fuerteconv}
Una función diferenciable $\mathcal{L}:\mathbb{R}^p\to\mathbb{R}$ es
\textbf{$\mu$-fuertemente convexa} con constante $\mu > 0$ si para todo
$\bm{\theta}, \bm{\phi} \in \mathbb{R}^p$:
\[
  \mathcal{L}(\bm{\phi}) \geq \mathcal{L}(\bm{\theta})
  + \langle\nabla\mathcal{L}(\bm{\theta}),\bm{\phi}-\bm{\theta}\rangle
  + \frac{\mu}{2}\|\bm{\phi}-\bm{\theta}\|_2^2.
\]
\end{definicion}

La convexidad fuerte es la convexidad ordinaria con un término cuadrático extra:
implica que la función crece al menos cuadráticamente en todas las direcciones, lo
que garantiza un mínimo global único y bien separado.

\begin{definicion}{Gradiente Lipschitz}
\label{def:lipschitz}
El gradiente de $\mathcal{L}$ es \textbf{$L$-Lipschitz} si para todo
$\bm{\theta}, \bm{\phi} \in \mathbb{R}^p$:
\[
  \|\nabla\mathcal{L}(\bm{\theta}) - \nabla\mathcal{L}(\bm{\phi})\|_2
  \leq L\,\|\bm{\theta}-\bm{\phi}\|_2.
\]
Para $\mathcal{L} \in C^2$, esta condición equivale a $\|\bm{H}(\bm{\theta})\|_2
\leq L$ para todo $\bm{\theta}$: el hessiano tiene autovalores acotados por $L$.
\end{definicion}

\begin{proposicion}{Convergencia del descenso de gradiente}
\label{prop:conv_gd}
Sea $\mathcal{L}$ $\mu$-fuertemente convexa con gradiente $L$-Lipschitz y mínimo
$\bm{\theta}^*$. El descenso de gradiente~\eqref{eq:gd} con tasa $\eta = 1/L$
satisface:
\[
  \mathcal{L}(\bm{\theta}^{(k)}) - \mathcal{L}(\bm{\theta}^*)
  \leq \left(1 - \frac{\mu}{L}\right)^k
  \bigl[\mathcal{L}(\bm{\theta}^{(0)}) - \mathcal{L}(\bm{\theta}^*)\bigr].
\]
\end{proposicion}

\begin{proof}
De la condición $L$-Lipschitz del gradiente se obtiene la cota de descenso:
\[
  \mathcal{L}(\bm{\theta}^{(k+1)}) \leq \mathcal{L}(\bm{\theta}^{(k)})
  - \frac{1}{2L}\|\nabla\mathcal{L}(\bm{\theta}^{(k)})\|_2^2.
\]
De la $\mu$-convexidad fuerte, $\|\nabla\mathcal{L}(\bm{\theta})\|_2^2 \geq
2\mu[\mathcal{L}(\bm{\theta})-\mathcal{L}(\bm{\theta}^*)]$. Combinando:
\[
  \mathcal{L}(\bm{\theta}^{(k+1)}) - \mathcal{L}(\bm{\theta}^*)
  \leq \left(1-\frac{\mu}{L}\right)
  \bigl[\mathcal{L}(\bm{\theta}^{(k)})-\mathcal{L}(\bm{\theta}^*)\bigr].
\]
Iterando $k$ veces se obtiene el resultado. $\square$
\end{proof}

El cociente $L/\mu$ es el \textbf{número de condición} del problema de optimización:
coincide, para funciones cuadráticas, con el número de condicionamiento del hessiano
$\kappa(\bm{H}) = \mu_{\max}/\mu_{\min}$ del Capítulo~\ref{ch:cap3}. Cuanto mayor
$L/\mu$, más lenta la convergencia: se necesitan $O((L/\mu)\log(1/\varepsilon))$
iteraciones para alcanzar precisión $\varepsilon$.

\subsection{Descenso de gradiente en lote completo y en mini-lotes}
\label{subsec:minibatch}

\dificultad{1}{2}

La pérdida empírica es $\mathcal{L}(\bm{\theta}) = \frac{1}{N}\sum_{i=1}^N
\ell_i(\bm{\theta})$ con $\ell_i(\bm{\theta}) = \ell(f(\bm{x}_i;\bm{\theta}),y_i)$.
El gradiente completo es $\nabla\mathcal{L} = \frac{1}{N}\sum_{i=1}^N\nabla\ell_i$:
requiere evaluar la red en las $N$ muestras, con costo $O(N)$ por iteración.
Para $N = 10^9$, una sola iteración puede tardar horas.

La solución es estimar el gradiente con un \textbf{mini-lote} aleatorio
$\mathcal{B} \subset \{1,\ldots,N\}$ de tamaño $B \ll N$:
\[
  \hat{\nabla}\mathcal{L}(\bm{\theta};\mathcal{B})
  = \frac{1}{B}\sum_{i\in\mathcal{B}} \nabla_{\bm{\theta}}\ell_i(\bm{\theta}).
\]
El costo por iteración es $O(B)$, independiente de $N$. El precio es que el
gradiente estimado tiene ruido. Analizar ese ruido es el objeto de la siguiente
sección.

%───────────────────────────────────────────────────────────────────────────────
\section{Descenso de gradiente estocástico y cálculo de Itô}
\label{sec:sgd_ito}
%───────────────────────────────────────────────────────────────────────────────

\dificultad{3}{4}

Esta sección desarrolla la conexión entre el SGD y las ecuaciones diferenciales
estocásticas. El resultado central ---que el SGD es la discretización de una SDE
de Langevin--- no solo explica el comportamiento del algoritmo sino que introduce
el lenguaje del cálculo estocástico que será imprescindible en el Capítulo~\ref{ch:cap12}
para los modelos de difusión.

\subsection{El gradiente estocástico como estimador con ruido}
\label{subsec:ruido_grad}

\begin{proposicion}{Insesgamiento del gradiente estocástico}
\label{prop:insesgado}
Sea $\mathcal{B}$ un mini-lote seleccionado uniformemente al azar de $\{1,\ldots,N\}$
con tamaño $|\mathcal{B}| = B$. Entonces:
\[
  \mathbb{E}_{\mathcal{B}}\bigl[\hat{\nabla}\mathcal{L}(\bm{\theta};\mathcal{B})\bigr]
  = \nabla\mathcal{L}(\bm{\theta}).
\]
\end{proposicion}

\begin{proof}
La esperanza sobre los mini-lotes aleatorios de tamaño $B$ seleccionados sin
reemplazamiento satisface:
$\mathbb{E}_{\mathcal{B}}\bigl[\frac{1}{B}\sum_{i\in\mathcal{B}}\nabla\ell_i\bigr]
= \frac{1}{N}\sum_{i=1}^N\nabla\ell_i = \nabla\mathcal{L}$,
donde se ha usado que cada índice $i$ aparece en $\binom{N-1}{B-1}/\binom{N}{B} =
B/N$ de los posibles mini-lotes. $\square$
\end{proof}

El error de estimación es:
\[
  \bm{\xi}(\bm{\theta};\mathcal{B})
  = \hat{\nabla}\mathcal{L}(\bm{\theta};\mathcal{B}) - \nabla\mathcal{L}(\bm{\theta}),
\]
con $\mathbb{E}[\bm{\xi}] = \bm{0}$ y covarianza:
\[
  \bm{C}(\bm{\theta}) = \mathbb{E}\bigl[\bm{\xi}\bm{\xi}^\top\bigr]
  = \frac{1}{B}\,\mathrm{Cov}_i[\nabla\ell_i(\bm{\theta})],
\]
donde $\mathrm{Cov}_i[\nabla\ell_i] = \frac{1}{N}\sum_{i=1}^N
(\nabla\ell_i-\nabla\mathcal{L})(\nabla\ell_i-\nabla\mathcal{L})^\top$ es la
covarianza empírica de los gradientes individuales. La varianza del estimador
escala como $1/B$: doblar el tamaño del mini-lote reduce el ruido a la mitad.

La regla de actualización del SGD es:
\begin{equation}
  \label{eq:sgd}
  \bm{\theta}^{(k+1)} = \bm{\theta}^{(k)}
  - \eta\,\nabla\mathcal{L}(\bm{\theta}^{(k)})
  - \eta\,\bm{\xi}^{(k)},
\end{equation}
donde $\bm{\xi}^{(k)}$ es el ruido del mini-lote en la iteración $k$.

\subsection{Elementos de cálculo de Itô}
\label{subsec:ito}

Para entender la dinámica de largo plazo del SGD se necesita el lenguaje de las
ecuaciones diferenciales estocásticas. Esta subsección introduce los conceptos
mínimos del cálculo de Itô de forma autocontenida.

\subsubsection*{El movimiento browniano}

\begin{definicion}{Movimiento browniano estándar}
\label{def:browniano}
Un \textbf{movimiento browniano} (o proceso de Wiener) en $\mathbb{R}^p$ es
un proceso estocástico $\bm{W}_t = (W_t^1,\ldots,W_t^p)$ tal que:
\begin{enumerate}
  \item $\bm{W}_0 = \bm{0}$ casi seguramente.
  \item Para $0 \leq s < t$, el incremento $\bm{W}_t - \bm{W}_s$ tiene
        distribución $\mathcal{N}(\bm{0},(t-s)\bm{I})$.
  \item Los incrementos en intervalos disjuntos son independientes.
  \item Las trayectorias $t \mapsto \bm{W}_t$ son continuas casi seguramente.
\end{enumerate}
\end{definicion}

El movimiento browniano es el límite de la suma de pasos aleatorios gaussianos
independientes cuando el paso tiende a cero. Sus trayectorias son continuas pero
no diferenciables en ningún punto: $\mathbb{E}[|W_{t+h}-W_t|^2] = h$, de modo
que el cociente $(W_{t+h}-W_t)/h$ tiene varianza $1/h \to \infty$ cuando $h\to 0$.
Por esta razón el movimiento browniano no admite derivada en el sentido ordinario,
y se necesita el cálculo de Itô para trabajar con él.

\subsubsection*{La integral de Itô}

La integral $\int_0^T f(t)\,dW_t$ no puede definirse como una integral de
Riemann-Stieltjes ordinaria porque las trayectorias de $W_t$ tienen variación
cuadrática infinita. La integral de Itô se define como límite en media cuadrática
de sumas de Riemann con el integrando evaluado en el extremo izquierdo de cada
subintervalo:
\[
  \int_0^T f(t)\,dW_t
  = \lim_{n\to\infty} \sum_{k=0}^{n-1} f(t_k)\,(W_{t_{k+1}}-W_{t_k}),
\]
donde $0 = t_0 < t_1 < \cdots < t_n = T$ es una partición. Esta elección del
punto de evaluación (extremo izquierdo) define el cálculo de Itô y es la que
aparece naturalmente en aplicaciones físicas y en el análisis del SGD.

La propiedad fundamental que distingue al cálculo de Itô del ordinario es la
\textbf{variación cuadrática}: el producto $(dW_t)^2$ no es negligible de segundo
orden sino que se comporta como $dt$:
\begin{equation}
  \label{eq:ito_rule}
  (dW_t^i)^2 = dt, \qquad dW_t^i\,dW_t^j = 0 \text{ para } i\neq j, \qquad
  dW_t^i\,dt = 0.
\end{equation}
Esta tabla de multiplicación (la \textbf{tabla de Itô}) es el corazón del cálculo
estocástico.

\subsubsection*{La fórmula de Itô}

La fórmula de Itô es el análogo estocástico de la regla de la cadena.

\begin{teorema}{Fórmula de Itô (caso unidimensional)}
\label{thm:ito}
Sea $X_t$ un proceso que satisface la SDE $dX_t = \mu_t\,dt + \sigma_t\,dW_t$,
y sea $g(t,x) \in C^{1,2}$. Entonces:
\[
  dg(t,X_t) = \frac{\partial g}{\partial t}\,dt
  + \frac{\partial g}{\partial x}\,dX_t
  + \frac{1}{2}\frac{\partial^2 g}{\partial x^2}\,(dX_t)^2.
\]
Usando $(dX_t)^2 = \sigma_t^2\,dt$ (de la tabla de Itô):
\[
  dg(t,X_t) = \left(\frac{\partial g}{\partial t}
  + \mu_t\frac{\partial g}{\partial x}
  + \frac{\sigma_t^2}{2}\frac{\partial^2 g}{\partial x^2}\right)dt
  + \sigma_t\frac{\partial g}{\partial x}\,dW_t.
\]
\end{teorema}

La diferencia con la regla de la cadena ordinaria es el término
$\frac{\sigma_t^2}{2}\frac{\partial^2 g}{\partial x^2}$: proviene de la variación
cuadrática del movimiento browniano y es de primer orden en $dt$ (no de segundo
orden como sería en cálculo ordinario). Este término es el responsable de toda la
riqueza del cálculo estocástico.

\subsection{La SDE de Langevin del SGD}
\label{subsec:langevin}

La actualización del SGD~\eqref{eq:sgd} puede reescribirse como:
\[
  \bm{\theta}^{(k+1)} - \bm{\theta}^{(k)}
  = -\eta\,\nabla\mathcal{L}(\bm{\theta}^{(k)}) - \eta\,\bm{\xi}^{(k)}.
\]
Si $\bm{\xi}^{(k)}$ es aproximadamente gaussiano con covarianza $\bm{C}(\bm{\theta}^{(k)})$
(lo que se justifica por el teorema central del límite cuando $B$ es moderado),
entonces $\eta\bm{\xi}^{(k)} \approx \sqrt{\eta}\,\bm{C}^{1/2}(\bm{\theta}^{(k)})
\,\Delta\bm{W}^{(k)}$, donde $\Delta\bm{W}^{(k)} \sim \mathcal{N}(\bm{0},\bm{I})$
son incrementos gaussianos independientes.

En el límite de paso pequeño $\eta \to 0$, manteniendo $\eta k = t$ fijo, el
SGD converge a la SDE de Itô:
\begin{equation}
  \label{eq:langevin_sgd}
  \boxed{d\bm{\theta}_t
  = -\nabla\mathcal{L}(\bm{\theta}_t)\,dt
  + \sqrt{\eta}\,\bm{C}^{1/2}(\bm{\theta}_t)\,d\bm{W}_t.}
\end{equation}
Esta es la \textbf{ecuación de Langevin sobremortigada} con potencial
$\mathcal{L}(\bm{\theta})$, temperatura efectiva $\eta$ y ruido multiplicativo
con covarianza $\bm{C}(\bm{\theta})$. Para $\bm{C} = \bm{I}$ (ruido isótropo)
es exactamente la dinámica de Langevin estándar de la mecánica estadística.

\begin{fisicaconexion}{La ecuación de Langevin en física estadística}
La ecuación~\eqref{eq:langevin_sgd} es la \textbf{ecuación de Langevin
sobremortigada}, que describe el movimiento de una partícula browniana en un
potencial $\mathcal{L}$ en el régimen de alta viscosidad. En física estadística
es el modelo estándar para la dinámica de una molécula en un solvente viscoso:
la fuerza $-\nabla\mathcal{L}$ es la fuerza conservativa del potencial, el término
estocástico $\sqrt{2k_BT}\,d\bm{W}_t$ es la fuerza de fluctuación del solvente
(ruido térmico), y el cociente $k_BT$ es la temperatura en unidades de energía.
En el SGD, la temperatura la controla el usuario mediante la tasa de aprendizaje
$\eta$ y el tamaño de mini-lote $B$ (que determina $\bm{C}$). El teorema de
fluctuación-disipación de Einstein establece que el coeficiente de difusión es
$D = k_BT/\gamma$: en el SGD, el análogo es que la ``difusión'' en el espacio de
parámetros escala como $\eta c$. Los cristales de proteínas se estudian con la
ecuación de Langevin a temperatura variable para encontrar el estado de mínima
energía libre: exactamente el mismo protocolo de recocido que usan los
calendarios de tasa de aprendizaje de la Sección~\ref{sec:lr_schedule}.
\end{fisicaconexion}

\begin{observacion}{Significado físico de cada término}
El término determinista $-\nabla\mathcal{L}\,dt$ produce descenso hacia el mínimo,
exactamente como el flujo de gradiente. El término estocástico
$\sqrt{\eta}\,\bm{C}^{1/2}\,d\bm{W}_t$ actúa como ruido térmico: hace fluctuar
los parámetros alrededor de la trayectoria determinista. La amplitud del ruido
escala como $\sqrt{\eta}$: una tasa de aprendizaje grande produce fluctuaciones
grandes. Esta es la razón por la que se usa una tasa de aprendizaje decreciente
durante el entrenamiento: al inicio conviene explorar (temperatura alta), al final
converger (temperatura baja).
\end{observacion}

\subsection{La ecuación de Fokker-Planck del SGD}
\label{subsec:fp_sgd}

La SDE~\eqref{eq:langevin_sgd} describe una trayectoria aleatoria en el espacio
de parámetros. La distribución de probabilidad $p(\bm{\theta},t)$ de encontrar
los parámetros en $\bm{\theta}$ en el tiempo $t$ satisface una ecuación
determinista: la ecuación de Fokker-Planck.

\begin{teorema}{Ecuación de Fokker-Planck del SGD}
\label{thm:fp_sgd}
Sea $p(\bm{\theta},t)$ la densidad de probabilidad de los parámetros bajo la
dinámica~\eqref{eq:langevin_sgd}. Entonces $p$ satisface:
\begin{equation}
  \label{eq:fp_sgd}
  \frac{\partial p}{\partial t}
  = \nabla_{\bm{\theta}} \cdot \bigl(p\,\nabla_{\bm{\theta}}\mathcal{L}\bigr)
  + \frac{\eta}{2}\,\nabla_{\bm{\theta}}^2 : \bigl(\bm{C}(\bm{\theta})\,p\bigr),
\end{equation}
donde $\nabla_{\bm{\theta}}^2 : \bm{A} = \sum_{i,j}\partial_{\theta_i\theta_j}A_{ij}$
es el operador de diferenciación de segundo orden contratado con $\bm{A}$.
\end{teorema}

\begin{proof}
Aplicamos la fórmula de Itô multidimensional a la función de prueba $g(\bm{\theta})$
y calculamos $\mathbb{E}[g(\bm{\theta}_t)]$. Para la SDE general
$d\bm{\theta}_t = \bm{f}(\bm{\theta}_t)\,dt + \bm{G}(\bm{\theta}_t)\,d\bm{W}_t$,
la fórmula de Itô da:
\[
  dg(\bm{\theta}_t)
  = \langle\nabla g, \bm{f}\rangle\,dt
  + \frac{1}{2}\mathrm{tr}\bigl(\bm{G}^\top(\nabla^2 g)\bm{G}\bigr)\,dt
  + \langle\nabla g, \bm{G}\,d\bm{W}_t\rangle.
\]
Para la SDE del SGD, $\bm{f} = -\nabla\mathcal{L}$ y $\bm{G} = \sqrt{\eta}\bm{C}^{1/2}$,
de modo que $\bm{G}\bm{G}^\top = \eta\bm{C}$. Tomando esperanza (el término con
$d\bm{W}_t$ tiene esperanza cero):
\[
  \frac{d}{dt}\mathbb{E}[g(\bm{\theta}_t)]
  = \mathbb{E}\!\left[
    -\langle\nabla g, \nabla\mathcal{L}\rangle
    + \frac{\eta}{2}\mathrm{tr}\bigl(\bm{C}(\bm{\theta}_t)\nabla^2 g\bigr)
  \right].
\]
Identificando $\frac{d}{dt}\mathbb{E}[g] = \int g\,\partial_t p\,d\bm{\theta}$
e integrando por partes (asumiendo que $p$ decae suficientemente rápido):
\begin{align*}
  \int g\,\partial_t p\,d\bm{\theta}
  &= \int \nabla g \cdot (p\nabla\mathcal{L})\,d\bm{\theta}
    \cdot(-1)^1
    + \frac{\eta}{2}\int g\,\nabla^2:(\bm{C}p)\,d\bm{\theta} \\
  &= \int g\left[
    \nabla\cdot(p\nabla\mathcal{L})
    + \frac{\eta}{2}\nabla^2:(\bm{C}p)
  \right]d\bm{\theta}.
\end{align*}
Como esto vale para toda función de prueba $g$, el integrando de ambos lados
debe coincidir, dando~\eqref{eq:fp_sgd}. $\square$
\end{proof}

\subsection{La distribución estacionaria: temperatura y recocido}
\label{subsec:estacionaria}

La ecuación de Fokker-Planck~\eqref{eq:fp_sgd} describe la evolución temporal
de la distribución de los parámetros. Su solución estacionaria
$\partial_t p^* = 0$ describe el régimen de largo plazo del SGD.

Para el caso simplificado de ruido isótropo $\bm{C}(\bm{\theta}) = c\bm{I}$
(covarianza constante), la ecuación estacionaria es:
\[
  0 = \nabla_{\bm{\theta}}\cdot\bigl(p^*\nabla\mathcal{L}\bigr)
  + \frac{\eta c}{2}\nabla_{\bm{\theta}}^2 p^*.
\]
Se propone la solución de tipo Gibbs: $p^*(\bm{\theta}) = Z^{-1}e^{-\mathcal{L}(\bm{\theta})/T_{\mathrm{eff}}}$.
Calculando:
\begin{align*}
  \nabla p^* &= -\frac{1}{T_{\mathrm{eff}}}p^*\nabla\mathcal{L}, \\
  \nabla\cdot(p^*\nabla\mathcal{L})
  &= \langle\nabla p^*, \nabla\mathcal{L}\rangle + p^*\nabla^2\mathcal{L}
   = -\frac{1}{T_{\mathrm{eff}}}p^*\|\nabla\mathcal{L}\|^2 + p^*\nabla^2\mathcal{L}, \\
  \nabla^2 p^*
  &= p^*\left(\frac{1}{T_{\mathrm{eff}}^2}\|\nabla\mathcal{L}\|^2
    - \frac{1}{T_{\mathrm{eff}}}\nabla^2\mathcal{L}\right).
\end{align*}
Sustituyendo en la ecuación estacionaria:
\[
  0 = p^*\left[
  -\frac{1}{T_{\mathrm{eff}}}\|\nabla\mathcal{L}\|^2 + \nabla^2\mathcal{L}
  + \frac{\eta c}{2}\left(\frac{\|\nabla\mathcal{L}\|^2}{T_{\mathrm{eff}}^2}
  - \frac{\nabla^2\mathcal{L}}{T_{\mathrm{eff}}}\right)
  \right].
\]
Para que esto sea cero para toda $\mathcal{L}$, se necesita $T_{\mathrm{eff}} =
\eta c/2$, con lo que los términos en $\|\nabla\mathcal{L}\|^2$ se cancelan.
La distribución estacionaria es entonces:

\begin{equation}
  \label{eq:gibbs}
  \boxed{p^*(\bm{\theta}) \propto \exp\!\left(-\frac{2\mathcal{L}(\bm{\theta})}{\eta c}\right).}
\end{equation}

\begin{observacion}{La distribución estacionaria como distribución de Gibbs}
\label{obs:gibbs}
La distribución~\eqref{eq:gibbs} es exactamente la distribución de Gibbs
(distribución de Boltzmann) de la mecánica estadística con potencial
$\mathcal{L}(\bm{\theta})$ y temperatura efectiva $T_{\mathrm{eff}} = \eta c/2$.
La masa de probabilidad se concentra en los mínimos de $\mathcal{L}$ (zonas de
baja energía), pero hay fluctuaciones alrededor de ellos con amplitud proporcional
a $T_{\mathrm{eff}}$.

Tres consecuencias prácticas inmediatas:

\begin{enumerate}
  \item \textbf{La tasa de aprendizaje como temperatura.} Una tasa de aprendizaje
        grande $\eta$ corresponde a temperatura alta: la distribución estacionaria
        es plana y los parámetros pueden escapar de mínimos locales poco profundos.
        Una tasa pequeña concentra la distribución en los mínimos más profundos.

  \item \textbf{El recocido como descenso de temperatura.} Reducir $\eta$ durante
        el entrenamiento es exactamente el algoritmo de recocido simulado
        (\textit{simulated annealing}) de la optimización combinatoria: se empieza
        con alta temperatura para explorar el espacio y se baja progresivamente para
        converger al mínimo. Los esquemas de tasa de aprendizaje de la
        Sección~\ref{sec:lr_schedule} son protocolos de recocido.

  \item \textbf{La covarianza del mini-lote como fuente de ruido controlable.}
        Aumentar $B$ reduce $c = \|\bm{C}\|/p$ (la varianza media por parámetro)
        y baja la temperatura efectiva. Para mantener la misma temperatura con
        el doble de tamaño de batch, se debe doblar $\eta$: esta es la
        \textbf{regla de escala lineal} del learning rate.
\end{enumerate}
\end{observacion}

\subsection{Épocas, iteraciones y el tamaño del mini-lote}
\label{subsec:epocas}

\begin{definicion}{Época e iteración}
\label{def:epoca}
Una \textbf{iteración} (o \textit{step}) es una actualización de parámetros
usando un mini-lote $\mathcal{B}$. Una \textbf{época} es una pasada completa sobre
el conjunto de entrenamiento, que consta de $\lceil N/B \rceil$ iteraciones.
\end{definicion}

La elección de $B$ involucra tres compromisos simultáneos. Desde el punto de vista
estadístico, un $B$ más grande reduce la varianza del gradiente como $1/B$. Desde
el punto de vista computacional, batches más grandes explotan mejor el paralelismo
de la GPU pero tienen rendimientos decrecientes: un batch de $2B$ no es el doble
de rápido que uno de $B$ en hardware moderno (la GPU ya está saturada para $B$
suficientemente grande). Desde el punto de vista de optimización, el ruido del SGD
tiene efectos de regularización implícita (la distribución estacionaria~\eqref{eq:gibbs}
penaliza mínimos estrechos en favor de mínimos amplios). En la práctica, $B \in
[32, 512]$ cubre la mayoría de las aplicaciones con hardware GPU estándar.

%───────────────────────────────────────────────────────────────────────────────
\section{SGD con momento}
\label{sec:momentum}
%───────────────────────────────────────────────────────────────────────────────

\dificultad{2}{3}

\subsection{El problema del zigzag y la motivación del momento}
\label{subsec:zigzag}

La Sección~\ref{subsec:conv_fuerte} estableció que la convergencia del descenso
de gradiente es geométrica con tasa $(1-\mu/L)^k$. Para un hessiano mal
condicionado ($\kappa = L/\mu \gg 1$), esta tasa es $(1-1/\kappa)^k \approx
e^{-k/\kappa}$: exponencialmente lenta. La causa geométrica se describió en la
Sección~\ref{sec:condicionamiento}: el gradiente oscila transversalmente en el
valle elongado mientras avanza lentamente hacia el fondo. El momento añade inercia
para amortiguar estas oscilaciones.

\subsection{SGD con momento de Polyak}
\label{subsec:polyak}

\begin{definicion}{SGD con momento}
\label{def:sgd_momentum}
El \textbf{SGD con momento} (Polyak, 1964) mantiene un vector de momento
$\bm{m}^{(k)}$ y actualiza:
\begin{align}
  \bm{m}^{(k+1)} &= \beta\bm{m}^{(k)} + (1-\beta)\hat{\nabla}\mathcal{L}(\bm{\theta}^{(k)}),
  \label{eq:momento} \\
  \bm{\theta}^{(k+1)} &= \bm{\theta}^{(k)} - \eta\,\bm{m}^{(k+1)},
  \label{eq:sgd_mom}
\end{align}
con $\beta \in (0,1)$ el \textbf{coeficiente de momento} (típicamente $\beta = 0.9$).
\end{definicion}

El vector $\bm{m}^{(k)}$ es la \textbf{media exponencialmente ponderada} de los
gradientes pasados:
\[
  \bm{m}^{(k)} = (1-\beta)\sum_{j=0}^k \beta^{k-j}\hat{\nabla}\mathcal{L}(\bm{\theta}^{(j)}).
\]
Gradientes de hace $k$ pasos atrás tienen peso $(1-\beta)\beta^k$: decaen
exponencialmente. El tiempo efectivo de memoria es $1/(1-\beta)$ pasos: para
$\beta = 0.9$, los gradientes de las últimas $\sim 10$ iteraciones tienen peso
significativo.

La ODE continua correspondiente a~\eqref{eq:momento}--\eqref{eq:sgd_mom} en el
límite $\eta,1-\beta\to 0$ es:
\[
  \ddot{\bm{\theta}} + (1-\beta)\dot{\bm{\theta}} + \nabla\mathcal{L}(\bm{\theta}) = \bm{0},
\]
que es exactamente la ecuación del \textbf{oscilador amortiguado} con constante
de amortiguamiento $1-\beta$ y potencial $\mathcal{L}$. Para $\beta$ cercano a $1$
(amortiguamiento débil), el sistema oscila pero converge más rápido que el
sobre-amortiguado (GD sin momento): la tasa óptima de convergencia mejora de
$O(1-1/\kappa)$ a $O(1-1/\sqrt{\kappa})$ (el número de iteraciones necesarias
escala como $\sqrt{\kappa}$ en vez de $\kappa$).

\begin{fisicaconexion}{Oscilador amortiguado y convergencia óptima}
La ODE $\ddot{\bm{\theta}} + (1-\beta)\dot{\bm{\theta}} + \nabla\mathcal{L}(\bm{\theta}) = \bm{0}$
es exactamente la ecuación del \textbf{oscilador amortiguado} de la mecánica
clásica. Para una función cuadrática $\mathcal{L} = \frac{\mu}{2}\theta^2$, se
tiene $\ddot{\theta} + 2\gamma\dot{\theta} + \omega_0^2\theta = 0$ con
$\omega_0^2 = \mu$ y $\gamma = (1-\beta)/2$. El discriminante $\Delta = \gamma^2 - \omega_0^2$
determina el régimen: sobre-amortiguado ($\Delta > 0$, análogo a GD sin momento),
críticamente amortiguado ($\Delta = 0$, convergencia más rápida posible), y
sub-amortiguado ($\Delta < 0$, oscilaciones que se extinguen). El régimen
críticamente amortiguado alcanza el equilibrio más rápido: exactamente el valor
óptimo de $\beta$ en el SGD con momento. En circuitos RLC, el mismo fenómeno
aparece en la respuesta al escalón: la descarga más rápida sin oscilaciones
ocurre para $R = 2\sqrt{L/C}$ (amortiguamiento crítico). El factor $\sqrt{\kappa}$
en la convergencia de Nesterov es el análogo de $1/\omega_0 = 1/\sqrt{\mu}$:
la frecuencia propia del sistema determina el tiempo de convergencia.
\end{fisicaconexion}

\subsection{Momento de Nesterov}
\label{subsec:nesterov}

\begin{definicion}{SGD con momento de Nesterov}
\label{def:nesterov}
El \textbf{momento de Nesterov} (NAG, Nesterov, 1983) evalúa el gradiente en
el punto anticipado:
\begin{align}
  \bm{m}^{(k+1)} &= \beta\bm{m}^{(k)} + \hat{\nabla}\mathcal{L}\bigl(\bm{\theta}^{(k)} - \eta\beta\bm{m}^{(k)}\bigr),
  \label{eq:nesterov_m}\\
  \bm{\theta}^{(k+1)} &= \bm{\theta}^{(k)} - \eta\,\bm{m}^{(k+1)}.
  \label{eq:nesterov_theta}
\end{align}
\end{definicion}

La diferencia con el momento de Polyak es sutil pero importante: en lugar de
calcular el gradiente en la posición actual $\bm{\theta}^{(k)}$, se calcula en
la posición anticipada $\bm{\theta}^{(k)} - \eta\beta\bm{m}^{(k)}$ (el punto
al que se llegaría si se aplicara el momento sin corrección). Esta anticipación
permite "frenar" antes de superar el mínimo.

\begin{proposicion}{Convergencia óptima de Nesterov}
\label{prop:nesterov_opt}
Para $\mathcal{L}$ convexa con gradiente $L$-Lipschitz, el SGD con momento de
Nesterov con tasa $\eta = 1/L$ satisface:
\[
  \mathcal{L}(\bm{\theta}^{(k)}) - \mathcal{L}(\bm{\theta}^*)
  \leq \frac{2L\|\bm{\theta}^{(0)}-\bm{\theta}^*\|_2^2}{k^2}.
\]
\end{proposicion}

La tasa $O(1/k^2)$ es la \textbf{tasa óptima de primer orden} para funciones
convexas: ningún algoritmo que use solo información de gradientes puede converger
más rápido en el peor caso (Nesterov, 1983). Para comparación, el GD sin momento
converge como $O(1/k)$.

%───────────────────────────────────────────────────────────────────────────────
\section{Optimizadores adaptativos}
\label{sec:adaptativos}
%───────────────────────────────────────────────────────────────────────────────

\dificultad{2}{2}

\subsection{Motivación: distintas direcciones necesitan distintas tasas}
\label{subsec:motivacion_adapt}

El hessiano mal condicionado tiene autovalores $\mu_1 \gg \mu_p > 0$. La tasa
de aprendizaje óptima en la dirección del autovector $\bm{v}_k$ es $\propto 1/\mu_k$:
las direcciones de alta curvatura necesitan pasos pequeños y las de baja curvatura
necesitan pasos grandes. Un único $\eta$ global es un compromiso subóptimo. La
solución ideal sería usar el paso $\bm{H}^{-1}$ (método de Newton), pero su costo
es $O(p^3)$. Los optimizadores adaptativos aproximan esta idea usando solo la
diagonal del hessiano, a costo $O(p)$.

\subsection{AdaGrad}
\label{subsec:adagrad}

\begin{definicion}{AdaGrad}
\label{def:adagrad}
AdaGrad (Duchi et al., 2011) acumula el cuadrado de los gradientes pasados
por parámetro:
\begin{align}
  G_j^{(k)} &= G_j^{(k-1)} + \bigl(\hat{\nabla}\mathcal{L}_j^{(k)}\bigr)^2,
  \label{eq:adagrad_G}\\
  \theta_j^{(k+1)} &= \theta_j^{(k)}
  - \frac{\eta}{\sqrt{G_j^{(k)}+\varepsilon}}\,\hat{\nabla}\mathcal{L}_j^{(k)},
  \label{eq:adagrad_update}
\end{align}
con $G_j^{(0)} = 0$ y $\varepsilon > 0$ una constante de estabilidad numérica.
\end{definicion}

El denominador $\sqrt{G_j^{(k)}}$ es la raíz cuadrática del promedio de los
cuadrados de los gradientes pasados del parámetro $j$: es una estimación de la
escala típica del gradiente en esa dirección. Dividiendo por esta escala, la tasa
efectiva para $\theta_j$ es $\eta/\sqrt{G_j^{(k)}}$: automáticamente pequeña para
parámetros con gradientes grandes y grande para parámetros con gradientes pequeños.

El problema de AdaGrad es que $G_j^{(k)}$ crece monótonamente: las tasas efectivas
decaen como $O(1/\sqrt{k})$ y tienden a cero en el largo plazo. Esto puede ser
deseable para problemas convexos (garantiza convergencia) pero es perjudicial para
problemas no convexos donde se necesita exploración continua.

\subsection{RMSProp}
\label{subsec:rmsprop}

\begin{definicion}{RMSProp}
\label{def:rmsprop}
RMSProp (Hinton, 2012) reemplaza la acumulación monótona de AdaGrad por una
media exponencialmente ponderada:
\begin{align}
  v_j^{(k)} &= \gamma\,v_j^{(k-1)} + (1-\gamma)\bigl(\hat{\nabla}\mathcal{L}_j^{(k)}\bigr)^2,
  \label{eq:rmsprop_v}\\
  \theta_j^{(k+1)} &= \theta_j^{(k)}
  - \frac{\eta}{\sqrt{v_j^{(k)}+\varepsilon}}\,\hat{\nabla}\mathcal{L}_j^{(k)},
  \label{eq:rmsprop_update}
\end{align}
con $\gamma \in (0,1)$ típicamente $0.9$ o $0.99$.
\end{definicion}

La media ponderada olvida gradientes antiguos con escala de tiempo $1/(1-\gamma)$:
la tasa efectiva se adapta a los cambios en la escala del gradiente a lo largo del
entrenamiento. RMSProp funciona bien en la práctica para redes recurrentes y fue
el optimizador preferido antes de Adam.

\subsection{Adam: momento de primer y segundo orden}
\label{subsec:adam}

\begin{definicion}{Adam}
\label{def:adam}
Adam (\textit{Adaptive Moment Estimation}, Kingma y Ba, 2015) combina el momento
de Polyak con la adaptación de RMSProp:
\begin{align}
  m_j^{(k)} &= \beta_1\,m_j^{(k-1)} + (1-\beta_1)\,\hat{\nabla}\mathcal{L}_j^{(k)},
  \label{eq:adam_m}\\
  v_j^{(k)} &= \beta_2\,v_j^{(k-1)} + (1-\beta_2)\,\bigl(\hat{\nabla}\mathcal{L}_j^{(k)}\bigr)^2,
  \label{eq:adam_v}\\
  \hat{m}_j^{(k)} &= \frac{m_j^{(k)}}{1-\beta_1^k}, \qquad
  \hat{v}_j^{(k)} = \frac{v_j^{(k)}}{1-\beta_2^k},
  \label{eq:adam_corr}\\
  \theta_j^{(k+1)} &= \theta_j^{(k)}
  - \frac{\eta\,\hat{m}_j^{(k)}}{\sqrt{\hat{v}_j^{(k)}}+\varepsilon}.
  \label{eq:adam_update}
\end{align}
Los valores estándar son $\beta_1 = 0.9$, $\beta_2 = 0.999$, $\varepsilon = 10^{-8}$
y $\eta = 3\times10^{-4}$.
\end{definicion}

Las correcciones~\eqref{eq:adam_corr} compensan el sesgo de inicialización: con
$m_j^{(0)} = v_j^{(0)} = 0$, los primeros valores de $m_j^{(k)}$ y $v_j^{(k)}$
subestimarían la magnitud real de los gradientes. Dividiendo por $1-\beta_1^k$ y
$1-\beta_2^k$ (que convergen a $1$ rápidamente) se obtienen estimadores centrados.

\subsection{Adam como precondicionamiento diagonal}
\label{subsec:adam_precon}

La actualización de Adam puede reescribirse como:
\[
  \bm{\theta}^{(k+1)} = \bm{\theta}^{(k)} - \eta\,\bm{D}^{(k)}\hat{\bm{m}}^{(k)},
\]
donde $\bm{D}^{(k)} = \mathrm{diag}\bigl(1/(\sqrt{\hat{v}_j^{(k)}}+\varepsilon)\bigr)$
es una matriz diagonal. La actualización es el gradiente suavizado $\hat{\bm{m}}$
multiplicado por el precondicionador diagonal $\bm{D}$.

El precondicionador $D_{jj} = 1/\sqrt{\hat{v}_j}$ escala cada componente por la
inversa de la desviación estándar estimada del gradiente en esa dirección. Para un
hessiano diagonal con entradas $H_{jj} = \mu_j$, la varianza del gradiente en la
dirección $j$ es $\sim \mu_j\sigma_\varepsilon^2$ (proporcional a la curvatura),
y $\sqrt{\hat{v}_j} \approx \mu_j^{1/2}\sigma_\varepsilon$, de modo que
$D_{jj} \approx 1/(\mu_j^{1/2}\sigma_\varepsilon)$. La actualización efectiva
escala como $\eta/\mu_j^{1/2}$: Adam es una aproximación al método de Newton
diagonal $\bm{H}^{-1}\nabla\mathcal{L}$ con precondicionar $\bm{H}^{-1/2}$
en lugar de $\bm{H}^{-1}$.

%───────────────────────────────────────────────────────────────────────────────
\section{Programación de la tasa de aprendizaje}
\label{sec:lr_schedule}
%───────────────────────────────────────────────────────────────────────────────

\dificultad{2}{3}

\subsection{La tasa como temperatura: justificación termodinámica}
\label{subsec:lr_temperatura}

La distribución estacionaria~\eqref{eq:gibbs} establece que el SGD en régimen
estacionario muestrea de $p^* \propto e^{-2\mathcal{L}/(\eta c)}$. Para
$\eta \to 0$, $p^*$ concentra toda su masa en el mínimo global (distribución de
Dirac): el SGD converge con certeza. Para $\eta$ grande, $p^*$ es aproximadamente
uniforme: el SGD explora el espacio. Un protocolo de entrenamiento eficiente
empieza con $\eta$ grande (exploración) y termina con $\eta$ pequeño (convergencia):
es el principio del recocido simulado.

\subsection{Calendarios de tasa de aprendizaje}
\label{subsec:calendarios}

Los calendarios más usados en la práctica son:

\medskip
\noindent\textbf{Decaimiento escalonado (\textit{step decay}).}
$\eta^{(k)} = \eta_0\,\gamma^{\lfloor k/s \rfloor}$ con $\gamma \in (0,1)$ y
periodo $s$. Sencillo de implementar y efectivo en problemas con estructura clara.

\medskip
\noindent\textbf{Decaimiento coseno (\textit{cosine annealing}).}
\[
  \eta^{(k)} = \eta_{\min} + \frac{1}{2}(\eta_{\max}-\eta_{\min})
  \left(1 + \cos\frac{\pi k}{K}\right),
\]
donde $K$ es el número total de iteraciones. La tasa cae suavemente de $\eta_{\max}$
a $\eta_{\min}$ siguiendo la curva de un coseno. Es el estándar en la mayoría de
los modelos de visión y física de altas energías (incluyendo CaloDiT).

\medskip
\noindent\textbf{Calentamiento lineal + decaimiento coseno (\textit{linear warmup}).}
Se incrementa $\eta$ linealmente desde $0$ hasta $\eta_{\max}$ durante las
primeras $k_w$ iteraciones, luego se aplica decaimiento coseno. El calentamiento
evita actualizaciones desestabilizadoras en las primeras iteraciones cuando los
parámetros están lejos de cualquier mínimo (los gradientes son grandes y Adam aún
no ha acumulado estadísticas de segundo momento).

\subsection{La regla de escala del learning rate con el batch size}
\label{subsec:lr_scale}

De la distribución estacionaria~\eqref{eq:gibbs}, la temperatura efectiva es
$T_{\mathrm{eff}} = \eta c/2$ donde $c \propto 1/B$ (la varianza del gradiente
estocástico). Para mantener la misma temperatura al multiplicar $B$ por $k$:
\[
  T_{\mathrm{eff}} = \frac{\eta c}{2} = \text{cte.}
  \implies \eta' = k\,\eta.
\]
Esta es la \textbf{regla de escala lineal} (Goyal et al., 2017): multiplicar el
batch size por $k$ y la tasa de aprendizaje por $k$ preserva la dinámica efectiva
del SGD. En la práctica, la regla lineal funciona bien para $k \leq 8$; para
escalas mayores, la regla de la raíz cuadrada $\eta' = \sqrt{k}\,\eta$ es más
estable.

%───────────────────────────────────────────────────────────────────────────────
\section{Inicialización de parámetros}
\label{sec:init}
%───────────────────────────────────────────────────────────────────────────────

\dificultad{2}{2}

\subsection{Por qué la inicialización importa}
\label{subsec:init_importa}

La inicialización $\bm{\theta}^{(0)}$ determina la cuenca de atracción a la que
el algoritmo converge. Pero hay un efecto más inmediato: una mala inicialización
puede causar activaciones saturadas (todas las neuronas sigmoide en la región
plana) o explosión de gradiente (varianza de activaciones que crece exponencialmente
con la profundidad), haciendo el entrenamiento imposible desde las primeras
iteraciones. La inicialización correcta mantiene la escala de las activaciones y
los gradientes aproximadamente constante a través de las capas.

El análisis es el siguiente: si los pesos de la capa $l$ son independientes con
media cero y varianza $\sigma_l^2$, y la entrada $\bm{a}^{(l-1)}$ tiene varianza
$V_{l-1}$, entonces la varianza de la pre-activación es:
\[
  \mathrm{Var}[z_j^{(l)}] = \mathrm{Var}\!\left[\sum_{k=1}^{n_{l-1}} W_{jk}^{(l)}a_k^{(l-1)}\right]
  = n_{l-1}\,\sigma_l^2\,V_{l-1},
\]
asumiendo que los pesos y las activaciones son independientes y las activaciones
tienen media cero. Para que $V_l = V_{l-1}$ (la varianza se preserva), se necesita
$n_{l-1}\,\sigma_l^2 = 1$, es decir $\sigma_l^2 = 1/n_{l-1}$.

\subsection{Inicialización de Xavier / Glorot}
\label{subsec:xavier}

Para activaciones con derivada $\approx 1$ en el origen (tanh linealizado,
sigmoide normalizada), el mismo análisis en el backward pass requiere
$\sigma_l^2 = 1/n_l$. El compromiso entre forward y backward da:

\begin{definicion}{Inicialización de Xavier}
\label{def:xavier}
La \textbf{inicialización de Xavier} (Glorot y Bengio, 2010) inicializa los pesos
de la capa $l$ como:
\[
  W_{jk}^{(l)} \sim \mathcal{U}\!\left(-\sqrt{\frac{6}{n_{l-1}+n_l}},\,\sqrt{\frac{6}{n_{l-1}+n_l}}\right)
  \quad\text{o equivalentemente}\quad
  W_{jk}^{(l)} \sim \mathcal{N}\!\left(0,\,\frac{2}{n_{l-1}+n_l}\right).
\]
\end{definicion}

Esta inicialización hace que la varianza de las activaciones sea aproximadamente
constante en el forward pass y que la varianza de los gradientes sea constante en
el backward pass.

\subsection{Inicialización de He / Kaiming}
\label{subsec:he}

Para ReLU, la mitad de las pre-activaciones son negativas y se mapean a cero.
La varianza de la post-activación es la mitad de la de la pre-activación:
$V_l = \frac{1}{2}n_{l-1}\sigma_l^2 V_{l-1}$. Para preservar la varianza se
necesita $\sigma_l^2 = 2/n_{l-1}$:

\begin{definicion}{Inicialización de He}
\label{def:he}
La \textbf{inicialización de He} (He et al., 2015) inicializa los pesos de la
capa $l$ con:
\[
  W_{jk}^{(l)} \sim \mathcal{N}\!\left(0,\,\frac{2}{n_{l-1}}\right).
\]
\end{definicion}

\begin{observacion}{Inicialización como condición inicial del integrador}
Desde la perspectiva del flujo de gradiente, la inicialización es la condición
inicial $\bm{\theta}^{(0)}$ del integrador. La inicialización de He garantiza
que el estado inicial del sistema (activaciones y gradientes) está en un régimen
de varianza $O(1)$: la "energía" del sistema es finita y el integrador puede dar
los primeros pasos de forma estable. Una mala inicialización es análoga a comenzar
una integración de Verlet con velocidades o aceleraciones infinitas: el integrador
diverge inmediatamente.
\end{observacion}

%───────────────────────────────────────────────────────────────────────────────
\section{Resumen del Capítulo 5}
\label{sec:resumen5}
%───────────────────────────────────────────────────────────────────────────────

El capítulo construyó una jerarquía de algoritmos de optimización, cada uno como
respuesta a una limitación del anterior.

El descenso de gradiente es la discretización de Euler de la ODE del flujo de
gradiente. Su convergencia es geométrica con tasa controlada por el número de
condicionamiento del hessiano. El SGD aproxima el gradiente con mini-lotes,
reduciendo el costo por iteración a $O(B)$ a costa de introducir ruido. Ese ruido
no es un defecto: convierte el flujo de gradiente en una SDE de Langevin cuya
distribución estacionaria es la distribución de Gibbs $p^* \propto e^{-2\mathcal{L}/\eta c}$.
La tasa de aprendizaje actúa como temperatura termodinámica, y su decaimiento
durante el entrenamiento es el protocolo de recocido que lleva al sistema del
régimen de exploración al de convergencia.

El momento añade inercia al integrador, mejorando la convergencia de $O(\kappa)$
a $O(\sqrt{\kappa})$ iteraciones. Adam combina momento de primer y segundo orden
para adaptar la tasa por parámetro, aproximando la acción del hessiano inverso
diagonal. La inicialización de He establece las condiciones iniciales que permiten
entrenar redes profundas con ReLU sin saturación ni explosión.

El capítulo deja una pregunta abierta: el gradiente $\nabla_{\bm{\theta}}\mathcal{L}$
se ha usado en todos los algoritmos, pero no se ha explicado cómo calcularlo
eficientemente para un MLP de $L$ capas y $p \sim 10^9$ parámetros. Hacerlo por
diferencias finitas requeriría $O(p)$ evaluaciones del forward pass por iteración:
computacionalmente imposible. El Capítulo~\ref{ch:cap6} presenta el algoritmo que
lo hace en $O(1)$ evaluaciones forward: la retropropagación, que es la aplicación
del método adjunto al grafo computacional del Capítulo~\ref{ch:cap4}.

%───────────────────────────────────────────────────────────────────────────────
\label{sec:estrella5}
\begin{estrella}{Optimización del punto de operación de un reactor con Adam}

\textbf{El problema.}
El modelo sustituto del Capítulo~\ref{ch:cap4} proporciona una función diferenciable
$f_{\mathrm{MLP}}(\bm{u}) \approx \bm{\Phi}(\bm{u})$ que mapea parámetros operacionales
$\bm{u} \in \mathbb{R}^d$ (posiciones de barras de control, caudal del refrigerante)
al flujo neutrónico $\bm{\Phi} \in \mathbb{R}^K$ en el núcleo. El control del reactor
requiere encontrar $\bm{u}^*$ tal que el flujo realice una distribución objetivo
$\bm{\Phi}_{\mathrm{obj}}$ (por ejemplo, distribución axial uniforme para minimizar
factores de potencia pico):
\[
  \bm{u}^* = \argmin_{\bm{u} \in \mathcal{U}}\;
  \mathcal{L}_{\mathrm{ctrl}}(\bm{u})
  = \tfrac{1}{2}\|f_{\mathrm{MLP}}(\bm{u}) - \bm{\Phi}_{\mathrm{obj}}\|_2^2,
\]
donde $\mathcal{U} \subset \mathbb{R}^d$ es el dominio de operación segura.

\textbf{Elección del optimizador.}
El paisaje de $\mathcal{L}_{\mathrm{ctrl}}$ hereda la no convexidad del sustituto.
Los distintos parámetros operacionales tienen escalas de sensibilidad muy diferentes:
la posición de una barra de control cambia el flujo local en un factor 10 para una
variación de pocos centímetros, mientras que la temperatura del refrigerante lo
cambia en menos de $1\%$ por kelvin. Adam es el optimizador natural para este
problema: la adaptación por parámetro normaliza automáticamente las distintas
escalas de sensibilidad, sin necesidad de escalar manualmente las variables.

\textbf{La tasa de aprendizaje como temperatura de control.}
La distribución estacionaria del SGD (ecuación~\eqref{eq:gibbs}) identifica $\eta$
con una temperatura que controla el nivel de exploración. En el contexto de control
del reactor, esto tiene una interpretación operativa: una tasa alta permite
explorar puntos de operación no estándar (posiblemente mejores) pero con mayor
incertidumbre; una tasa baja converge a un punto de operación preciso pero puede
quedar atrapada en un óptimo local subóptimo. El calentamiento lineal +
decaimiento coseno (Sección~\ref{sec:lr_schedule}) es el protocolo natural:
exploración inicial para mapear el espacio de operación, seguida de convergencia
fina hacia $\bm{u}^*$.

\textbf{Inicialización desde el punto de operación nominal.}
La inicialización de He es apropiada para entrenar el sustituto (pesos del MLP),
pero para la optimización de control la inicialización natural es $\bm{u}^{(0)} =
\bm{u}_{\mathrm{nom}}$ (punto de operación nominal): se sabe que está en la región
$\mathcal{U}$ y cerca de algún óptimo local razonable. El análisis de la
Sección~\ref{sec:init} aplica mutatis mutandis: la inicialización determina
la cuenca de atracción a la que el optimizador converge.

\textbf{Conexión con el capítulo.}
El gradiente $\nabla_{\bm{u}}\mathcal{L}_{\mathrm{ctrl}}$ se calcula por
backpropagation a través del sustituto (Capítulo~\ref{ch:cap6}): la composición
$f_{\mathrm{MLP}} \circ (\cdot)$ admite diferenciación automática exactamente
como el forward pass del MLP. Cada iteración de Adam cuesta un forward pass más
un backward pass sobre el sustituto: $O(p)$ operaciones, frente a las horas del
simulador de referencia. El proceso completo de optimización de control ---$10^3$
iteraciones de Adam sobre el sustituto--- es más barato que una sola evaluación
del simulador.
\end{estrella}

%───────────────────────────────────────────────────────────────────────────────
\section*{Ejercicios computacionales}
\addcontentsline{toc}{section}{Ejercicios computacionales}
%───────────────────────────────────────────────────────────────────────────────

\begin{ejerciciocomp}{Comparación de optimizadores sobre un paisaje de pérdida 2D}
\textbf{Objetivo:} visualizar las trayectorias de GD, SGD, SGD con momento y Adam
sobre un paisaje de pérdida bidimensional mal condicionado y verificar las
predicciones teóricas de convergencia.

\medskip\noindent\textbf{Función objetivo:}
\[
  \mathcal{L}(\theta_1, \theta_2) = \frac{1}{2}(\theta_1^2 + 100\theta_2^2),
  \quad \bm{\theta}^{(0)} = (1, 1)^\top, \quad \bm{\theta}^* = (0, 0)^\top.
\]

\medskip\noindent\textbf{Algoritmo:}
\begin{enumerate}
  \item Ejecutar $500$ iteraciones de GD, SGD (ruido $\sigma = 0.1$), SGD con momento
        ($\beta = 0.9$) y Adam con $\eta = 0.01$ para todos.
  \item Graficar las trayectorias $(\theta_1^{(k)}, \theta_2^{(k)})$ superpuestas
        sobre las curvas de nivel de $\mathcal{L}$.
  \item Trazar $\log\mathcal{L}(\bm{\theta}^{(k)})$ en función de $k$ para
        cada optimizador en una figura separada.
  \item Medir el número de iteraciones para alcanzar $\mathcal{L} < 10^{-4}$ y
        comparar con las predicciones teóricas: $O(\kappa)$ para GD, $O(\sqrt{\kappa})$
        para momento de Nesterov.
\end{enumerate}

\medskip\noindent\textbf{Verificación:} Adam debe converger en menos iteraciones
que GD para $\kappa = 100$, y las trayectorias de GD deben mostrar el patrón de
zigzag predicho por la Sección~\ref{subsec:zigzag}.
\end{ejerciciocomp}

\begin{ejerciciocomp}{Implementación del calendario de tasa de aprendizaje coseno}
\textbf{Objetivo:} implementar los tres calendarios de la Sección~\ref{subsec:calendarios}
y verificar el efecto de la temperatura sobre la distribución estacionaria del SGD.

\medskip\noindent\textbf{Algoritmo:}
\begin{enumerate}
  \item Implementar $\eta^{(k)}$ para decaimiento escalonado ($\gamma = 0.5$,
        $s = 100$), decaimiento coseno y calentamiento lineal ($k_w = 50$)
        + decaimiento coseno. Graficar los tres calendarios para $K = 500$ iteraciones.
  \item Minimizar $\mathcal{L}(\theta) = (\theta - 3)^2$ con SGD (ruido
        $\sigma = 1$) usando cada calendario. Graficar $\theta^{(k)}$ y
        $\mathcal{L}(\theta^{(k)})$ en función de $k$.
  \item Estimar empíricamente la distribución estacionaria tomando los últimos
        $100$ valores de $\theta^{(k)}$ para cada calendario. ¿Se parece a la
        distribución de Gibbs $p^* \propto e^{-2\mathcal{L}(\theta)/\eta_{\mathrm{final}} c}$?
  \item Para el calendario coseno, verificar la regla de escala lineal: comparar
        las trayectorias para $(B=32, \eta=0.01)$ y $(B=128, \eta=0.04)$. ¿Son
        estadísticamente equivalentes?
\end{enumerate}
\end{ejerciciocomp}

%───────────────────────────────────────────────────────────────────────────────
\begin{notasbib}
El descenso de gradiente en tiempo continuo (flujo de gradiente) fue estudiado
sistemáticamente en la teoría de control y optimización convexa desde la década
de 1960. La perspectiva del método de Euler como discretización de la ODE del
flujo de gradiente aparece en los libros clásicos de análisis numérico de ODEs,
pero su aplicación sistemática al análisis del aprendizaje profundo fue
desarrollada por Hardt, Recht y Singer (2016) y formalizada en la literatura
de física matemática del aprendizaje.

El análisis de convergencia para funciones fuertemente convexas con gradiente
Lipschitz es un resultado clásico de la teoría de optimización convexa,
presentado de forma canónica en Nesterov (2004)~\cite{nesterov2004introductory}
y Nocedal y Wright (2006)~\cite{nocedal2006numerical}. La mejora de la tasa de
convergencia de $O(\kappa)$ a $O(\sqrt{\kappa})$ mediante el momento acelerado fue
establecida por Nesterov (1983) como óptima en la clase de métodos de primer orden.

La conexión entre el SGD y la ecuación de Langevin fue explorada por Chaudhari
y Soatto (2018), Li et al.\ (2017), y formalizada con herramientas del cálculo
estocástico por varios grupos. La interpretación de la tasa de aprendizaje como
temperatura y la distribución estacionaria de Gibbs aparecen en Welling y Teh
(2011)~\cite{welling2011bayesian} en el contexto del SGD bayesiano (SGLD). La
regla de escala lineal del learning rate con el batch size fue establecida
empíricamente por Goyal et al.\ (2017)~\cite{goyal2017accurate} en el contexto
del entrenamiento distribuido de ResNets, y derivada teóricamente de la distribución
estacionaria por varios autores.

AdaGrad fue introducido por Duchi, Hazan y Singer (2011)~\cite{duchi2011adaptive},
que probaron convergencia $O(1/\sqrt{k})$ para funciones convexas no suaves.
RMSProp fue propuesto por Hinton en notas de curso no publicadas (2012) como
solución al problema de tasas decrecientes de AdaGrad. Adam fue publicado por
Kingma y Ba (2015)~\cite{kingma2015adam} y se convirtió rápidamente en el
optimizador estándar para la mayoría de las aplicaciones de aprendizaje profundo.

La inicialización de Xavier fue derivada por Glorot y Bengio (2010)~\cite{glorot2010understanding},
que identificaron la varianza de las activaciones como el factor clave para la
propagación estable del gradiente. La inicialización de He (He et al.,
2015)~\cite{he2015delving} extendió el análisis a activaciones ReLU y permitió
entrenar redes de cientos de capas por primera vez.

La conexión entre el SGD y el recocido simulado tiene raíces en Kirkpatrick,
Gelatt y Vecchi (1983), que introdujeron el recocido simulado para la
optimización combinatoria, y en la teoría clásica de la física estadística del
equilibrio y fuera del equilibrio. En el contexto de los modelos de difusión del
Capítulo~\ref{ch:cap12}, la ecuación de Langevin y su distribución estacionaria
aparecen de forma central en la derivación del proceso de inversión temporal.
\end{notasbib}

%───────────────────────────────────────────────────────────────────────────────
\section*{Ejercicios resueltos}
\addcontentsline{toc}{section}{Ejercicios resueltos}
%───────────────────────────────────────────────────────────────────────────────

\begin{ejercicio}
\label{ej:flujo_grad_cuad}
\textbf{(Flujo de gradiente para una función cuadrática).}

Sea $\mathcal{L}(\bm{\theta}) = \frac{1}{2}\bm{\theta}^\top\bm{H}\bm{\theta}$
con $\bm{H}$ simétrica definida positiva. Resolver la ODE del flujo de gradiente
y mostrar que la solución converge al mínimo $\bm{\theta}^* = \bm{0}$.

\medskip\noindent\textit{Solución.}
La ODE es $\dot{\bm{\theta}} = -\bm{H}\bm{\theta}$. En la base de autovectores
$\bm{H} = \bm{Q}\bm{\Lambda}\bm{Q}^\top$, con $\bm{\phi} = \bm{Q}^\top\bm{\theta}$:
$\dot{\phi}_k = -\mu_k\phi_k$. Las soluciones son $\phi_k(t) = \phi_k(0)e^{-\mu_k t}$.
Como $\mu_k > 0$, cada componente decae exponencialmente: $\bm{\theta}(t) =
\bm{Q}\mathrm{diag}(e^{-\mu_k t})\bm{Q}^\top\bm{\theta}(0) \to \bm{0}$ cuando
$t\to\infty$. La tasa de convergencia está dominada por el modo más lento
$e^{-\mu_{\min}t}$: las componentes en las direcciones de menor curvatura
tardan más en converger, exactamente como predice la Proposición~\ref{prop:conv_cuad}.
$\square$
\end{ejercicio}

\begin{ejercicio}
\label{ej:ito_calculo}
\textbf{(Aplicación de la fórmula de Itô).}

Sea $X_t$ un movimiento browniano estándar ($dX_t = dW_t$). Calcular $d(X_t^2)$
usando la fórmula de Itô y verificar que $\mathbb{E}[X_t^2] = t$.

\medskip\noindent\textit{Solución.}
Aplicando el Teorema~\ref{thm:ito} con $g(x) = x^2$:
\[
  d(X_t^2) = 2X_t\,dX_t + \frac{1}{2}\cdot 2\cdot (dX_t)^2
  = 2X_t\,dW_t + (dW_t)^2.
\]
Usando la tabla de Itô: $(dW_t)^2 = dt$. Entonces:
\[
  d(X_t^2) = 2X_t\,dW_t + dt.
\]
Integrando de $0$ a $t$ y tomando esperanza (la integral estocástica tiene
esperanza cero):
\[
  \mathbb{E}[X_t^2] = X_0^2 + t = t,
\]
ya que $X_0 = 0$. Nótese que la fórmula ordinaria $d(x^2) = 2x\,dx$ no
incluiría el término $dt$: ese término es la corrección de Itô, debida a la
variación cuadrática del movimiento browniano. Omitirlo daría $\mathbb{E}[X_t^2]
= 0$, que es incorrecto. $\square$
\end{ejercicio}

\begin{ejercicio}
\label{ej:adam_manual}
\textbf{(Dos iteraciones de Adam).}

Sea $\mathcal{L}(\theta) = (\theta-3)^2$ y $\theta^{(0)} = 0$. Realizar dos
iteraciones de Adam con $\eta = 0.1$, $\beta_1 = 0.9$, $\beta_2 = 0.999$,
$\varepsilon = 10^{-8}$, usando el gradiente exacto.

\medskip\noindent\textit{Solución.}
$\nabla\mathcal{L}(\theta) = 2(\theta-3)$.

\textit{Iteración $k=1$:} $g^{(1)} = 2(0-3) = -6$.
\begin{align*}
  m^{(1)} &= 0.9\cdot 0 + 0.1\cdot(-6) = -0.6, \\
  v^{(1)} &= 0.999\cdot 0 + 0.001\cdot 36 = 0.036, \\
  \hat{m}^{(1)} &= -0.6/(1-0.9) = -6, \\
  \hat{v}^{(1)} &= 0.036/(1-0.999) = 36, \\
  \theta^{(1)} &= 0 - 0.1\cdot(-6)/(\sqrt{36}+10^{-8}) = 0.1.
\end{align*}

\textit{Iteración $k=2$:} $g^{(2)} = 2(0.1-3) = -5.8$.
\begin{align*}
  m^{(2)} &= 0.9\cdot(-0.6) + 0.1\cdot(-5.8) = -0.54 - 0.58 = -1.12, \\
  v^{(2)} &= 0.999\cdot 0.036 + 0.001\cdot 33.64 = 0.035964 + 0.03364 = 0.069604, \\
  \hat{m}^{(2)} &= -1.12/(1-0.81) = -5.895, \\
  \hat{v}^{(2)} &= 0.069604/(1-0.998) = 34.802, \\
  \theta^{(2)} &= 0.1 - 0.1\cdot(-5.895)/(\sqrt{34.802}+10^{-8}) \approx 0.1 + 0.0999 \approx 0.2.
\end{align*}

Cada iteración avanza $\approx 0.1$ hacia el mínimo $\theta^* = 3$,
independientemente de la escala del gradiente: la normalización por
$\sqrt{\hat{v}}$ hace que la actualización efectiva sea siempre $\approx \eta$
en magnitud. Esto ilustra la propiedad de Adam de mantener la magnitud de las
actualizaciones estable, independientemente de la escala del problema. $\square$
\end{ejercicio}

%───────────────────────────────────────────────────────────────────────────────
\section*{Ejercicios propuestos}
\addcontentsline{toc}{section}{Ejercicios propuestos}
%───────────────────────────────────────────────────────────────────────────────

\begin{ejercicio}
\textbf{(Estabilidad del método de Euler para el flujo de gradiente).}
Sea $\mathcal{L}(\theta_1,\theta_2) = \frac{1}{2}(\theta_1^2 + 100\theta_2^2)$
($\kappa = 100$).
\begin{enumerate}[label=(\alph*)]
  \item Calcule la tasa de aprendizaje máxima $\eta_{\max}$ para la estabilidad
        de Euler. ¿Qué ocurre para $\eta > \eta_{\max}$?
  \item Con $\eta = \eta^* = 2/(\mu_1+\mu_2)$ (tasa óptima), calcule cuántas
        iteraciones se necesitan para que $\|\bm{\theta}^{(k)}-\bm{\theta}^*\|
        < 10^{-3}$ partiendo de $\bm{\theta}^{(0)} = (1,1)^\top$.
  \item Repita con momento de Polyak ($\beta = 0.9$). ¿Cuántas iteraciones
        se necesitan ahora?
  \item Compare con el resultado teórico $O(\sqrt{\kappa})$ para momento de
        Nesterov.
\end{enumerate}
\end{ejercicio}

\begin{ejercicio}
\textbf{(La fórmula de Itô para la pérdida a lo largo del SGD).}
Sea $\mathcal{L}(\bm{\theta})$ una función suave y $\bm{\theta}_t$ la solución
de la SDE de Langevin~\eqref{eq:langevin_sgd} con $\bm{C} = c\bm{I}$.
\begin{enumerate}[label=(\alph*)]
  \item Aplicar la fórmula de Itô multidimensional para calcular $d\mathcal{L}
        (\bm{\theta}_t)$.
  \item Mostrar que $\frac{d}{dt}\mathbb{E}[\mathcal{L}(\bm{\theta}_t)] =
        -\mathbb{E}[\|\nabla\mathcal{L}\|^2] + \frac{\eta c}{2}\mathbb{E}
        [\nabla^2\mathcal{L}(\bm{\theta}_t)]$.
  \item Interpretar el término $\frac{\eta c}{2}\mathbb{E}[\nabla^2\mathcal{L}]$:
        ¿qué ocurre en el mínimo $\bm{\theta}^*$ donde $\nabla\mathcal{L} = \bm{0}$
        pero $\nabla^2\mathcal{L} \succ 0$? ¿Por qué la pérdida esperada no
        converge a cero sino a un valor positivo?
  \item Para la función cuadrática $\mathcal{L} = \frac{1}{2}\bm{\theta}^\top
        \bm{H}\bm{\theta}$, resolver la ODE para $\mathbb{E}[\mathcal{L}]$ y
        calcular el valor estacionario.
\end{enumerate}
\end{ejercicio}

\begin{ejercicio}
\textbf{(Implementación comparativa de optimizadores).}
Minimizar $\mathcal{L}(\bm{\theta}) = \frac{1}{2}\bm{\theta}^\top\bm{H}\bm{\theta}
- \bm{b}^\top\bm{\theta}$ con $\bm{H} = \mathrm{diag}(1, 10, 100, 1000)$,
$\bm{b} = (1,1,1,1)^\top$.
\begin{enumerate}[label=(\alph*)]
  \item Implemente GD, SGD (con ruido gaussiano añadido), SGD con momento de
        Polyak, y Adam.
  \item Para cada optimizador, trace la curva de convergencia $\mathcal{L}
        (\bm{\theta}^{(k)})-\mathcal{L}^*$ en función de $k$ (en escala
        logarítmica).
  \item ¿Qué optimizador converge más rápido? ¿Es consistente con el análisis
        teórico?
  \item Añada ruido gaussiano al gradiente con covarianza $\sigma^2\bm{I}$
        para simular SGD. ¿Cómo cambia la convergencia? ¿El SGD converge
        a $\mathcal{L}^*$ o a un valor mayor?
\end{enumerate}
\end{ejercicio}

\begin{ejercicio}
\textbf{(La distribución estacionaria del SGD).}
Para la SDE $d\theta = -\theta\,dt + \sqrt{\eta}\,dW_t$ (gradiente de
$\mathcal{L} = \theta^2/2$) con temperatura efectiva $\eta$:
\begin{enumerate}[label=(\alph*)]
  \item Escribir la ecuación de Fokker-Planck correspondiente.
  \item Encontrar la distribución estacionaria $p^*(\theta)$ resolviendo
        $\partial_t p^* = 0$.
  \item Verificar que $p^*(\theta) \propto e^{-\theta^2/\eta}$
        (distribución de Gibbs con temperatura $\eta/2$).
  \item ¿Cuál es la varianza de $p^*$? ¿Cómo depende del mínimo al que
        converge el SGD (estrecho vs. ancho) de la temperatura $\eta$?
\end{enumerate}
\end{ejercicio}

\begin{ejercicio}
\textbf{(Inicialización y propagación de varianza).}
Sea un MLP con $L = 5$ capas de anchura constante $m = 128$, activación ReLU
e inicialización $W_{jk}^{(l)} \sim \mathcal{N}(0,\sigma^2)$.
\begin{enumerate}[label=(\alph*)]
  \item Para $\sigma^2 = 1/m$, calcule analíticamente la varianza de
        $\bm{z}^{(l)}$ en función de $l$. ¿Explota o se desvanece?
  \item Para $\sigma^2 = 2/m$ (inicialización de He), repita el cálculo.
        ¿Qué ocurre ahora?
  \item Verifique numéricamente propagando un vector de entrada gaussiano
        $\bm{x} \sim \mathcal{N}(\bm{0},\bm{I})$ y midiendo
        $\mathrm{Var}[\bm{z}^{(l)}]$ para $l = 1,\ldots,5$.
  \item ¿Qué inicialización se necesita para activación sigmoide (cuya
        derivada esperada es $\approx 1/4$)?
\end{enumerate}
\end{ejercicio}
